% ============================================================================
\section{Formulación del problema}

\subsection{Problema general}

¿En qué medida un modelo de computación de reservorio cuántico mejora la predicción de las revisiones posteriores de la primera estimación del PBI peruano, en comparación con modelos econométricos y de aprendizaje automático clásicos, cuando la evaluación respeta estrictamente la información disponible en cada fecha de publicación?

\subsection{Problemas específicos}

\begin{enumerate}
	\item ¿Cuáles fueron las propiedades estadísticas de las revisiones del PBI peruano en términos de sesgo medio, dispersión, persistencia y heterogeneidad temporal?
	\item ¿Qué capacidad predictiva presentaron los modelos econométricos y de aprendizaje automático clásicos para anticipar las revisiones posteriores?
	\item ¿QRC aportó una mejora predictiva fuera de muestra respecto del mejor benchmark clásico y de un reservorio clásico comparable?
	\item ¿Bajo qué combinaciones de tamaño muestral, persistencia, dimensionalidad y razón señal-ruido se observó una mejora reproducible de QRC en datos sintéticos calibrados?
	\item Si se encontrara una mejora predictiva robusta, ¿en qué medida podría modificar una evaluación económica basada en información disponible en tiempo real?
\end{enumerate}
